\chapter{Gauduchon metrics and linear mass}
\label{chap:gauduchon}

A positive current can have a conserved linear mass even when its Monge--Amp\`ere mass varies. Gauduchon metrics provide the test forms for this conservation law: their $(n-1)$st powers annihilate $\ddc u$ under integration. Varying the metric then gives a criterion for a smooth background to admit a positive current. The same conformal construction, applied to a smooth solution metric, produces the measure that characterizes the range of the linearized Monge--Amp\`ere operator. These two uses of Gauduchon metrics connect positivity of currents with solvability of scalar equations.

Assume $n\geq2$ throughout this chapter. In dimension one every Hermitian form is closed, and the corresponding integration identities follow directly from Stokes' theorem.

\section{Conformal Gauduchon metrics}

The condition needed for linear mass conservation can be achieved in every conformal class. Its construction reduces to finding a positive generator of the adjoint kernel of a scalar elliptic operator.

\begin{definition}\label{def:gauduchon-metric}
A Hermitian form $\gamma$ is \emph{Gauduchon} if
\[
 \ddc(\gamma^{n-1})=0.
\]
\end{definition}

\begin{theorem}[Gauduchon's theorem]\label{thm:gauduchon}
Every Hermitian form $\eta$ on a connected compact complex manifold of dimension at least two has a conformal Gauduchon representative. More precisely, there is a smooth strictly positive function $\rho$, unique up to multiplication by a positive constant, such that
\[
 \ddc(\rho\eta^{n-1})=0.
\]
The associated Gauduchon form is $\gamma=\rho^{1/(n-1)}\eta$.
\end{theorem}
\Needspace{8\baselineskip}
\begin{proof}
We use the standard elliptic Fredholm theorem, elliptic regularity and the strong maximum principle. For a smooth real function $v$ on $X$, define
\[
 Lv=n\frac{\ddc v\wedge\eta^{n-1}}{\eta^n}.
\]
This real second-order operator has positive definite principal symbol and no zeroth-order term. The maximum principle gives $\ker L=\R$. Its index is zero: deform its positive principal symbol to that of a Laplace operator, or use the elliptic path joining the two operators. Hence its formal adjoint $L^*$ with respect to $\eta^n$ has a one-dimensional kernel, generated by a nonzero smooth real function $\rho$.

We first prove that $\int_X\rho\eta^n\ne0$. Otherwise the Fredholm alternative would solve $Lv=1$, contrary to the maximum principle at a maximum of $v$. Rescale so that $\int_X\rho\eta^n=1$.

Suppose that $\rho$ is negative somewhere. Choose a nonzero smooth $h\geq0$ supported where $\rho<0$, and put $c=\int_Xh\rho\eta^n<0$. The function $h-c$ is strictly positive and orthogonal to the adjoint kernel. The Fredholm alternative would therefore solve $Lv=h-c>0$, again impossible. Thus $\rho\geq0$.

In real local coordinates $(x_1,\ldots,x_{2n})$, write $L^*=A^{ij}\partial_i\partial_j+B^i\partial_i+C$, where $\partial_i=\partial/\partial x_i$, the coefficients are smooth, and repeated indices are summed from $1$ to $2n$. Its leading symbol is positive definite. Since
\[
 (L^*-\max(C,0))\rho=-\max(C,0)\rho\leq0,
\]
the strong minimum principle for an operator with nonpositive zeroth-order term shows that a zero of $\rho$ would force $\rho$ to vanish identically on the connected manifold. Hence $\rho>0$.

Finally, integration by parts gives
\[
 \int_X\rho Lv\,\eta^n
 =n\int_Xv\,\ddc(\rho\eta^{n-1}).
\]
Thus $L^*\rho=0$ is precisely the required equation. The one-dimensional adjoint kernel proves uniqueness up to positive scaling.
\end{proof}

The K\"ahler condition $\mathrm{d}\gamma=0$ implies the balanced condition $\mathrm{d}(\gamma^{n-1})=0$, which in turn implies that $\gamma$ is Gauduchon. On surfaces the first two conditions coincide. The conformal theorem gives the Gauduchon condition; integration by parts against other mixed powers still requires their own differential identities.

\Needspace{12\baselineskip}
\section{Conservation of the trace mass}

Testing a positive current against the $(n-1)$st power of a Gauduchon metric gives a finite mass determined entirely by its smooth background. The identity follows directly from distributional integration by parts, so it applies to unbounded potentials as well.

\begin{proposition}\label{prop:gauduchon-trace}
Let $\theta$ be a smooth real $(1,1)$-form, let $u$ be a $\theta$-plurisubharmonic function, and let $\gamma$ be a Gauduchon metric. Then
\begin{equation}\label{eq:gauduchon-trace}
 \int_X(\theta+\ddc u)\wedge\gamma^{n-1}
 =\int_X\theta\wedge\gamma^{n-1}.
\end{equation}
The left side is the mass of a positive measure. It is finite even when $u$ is unbounded.
\end{proposition}
\begin{proof}
A quasi-plurisubharmonic function is locally integrable. The current $\ddc u$ therefore acts on the smooth test form $\gamma^{n-1}$ by
\[
 \langle\ddc u,\gamma^{n-1}\rangle
 =\langle u,\ddc(\gamma^{n-1})\rangle=0.
\]
This proves the identity. Positivity of $\theta+\ddc u$ and smooth positivity of $\gamma^{n-1}$ show that their product is a positive measure. Its mass is the finite number on the right.
\end{proof}

Here $\theta+\ddc u$ is the full distributional current, including any mass on the pole set of $u$. The non-pluripolar construction will discard this pole mass; \cref{ex:p1-raw-tail} illustrates the difference on a curve.

\begin{corollary}\label{cor:trace-bounds}
Fix a Hermitian form $\omega$. There are constants $0<a\leq b<+\infty$, depending only on $\omega$ and a chosen Gauduchon metric, such that every $u\in\PSH(X,\omega)$ satisfies
\[
 a\leq\int_X(\omega+\ddc u)\wedge\omega^{n-1}\leq b.
\]
\end{corollary}
\begin{proof}
Fix a Gauduchon form $\gamma$ on $X$ and choose positive constants $c,C$ with
$c\gamma^{n-1}\leq\omega^{n-1}\leq C\gamma^{n-1}$.
Wedge these inequalities with the positive current $\omega+\ddc u$ and integrate. By \cref{prop:gauduchon-trace}, the resulting bounds are
$cA$ and $CA$, where $A=\int_X\omega\wedge\gamma^{n-1}>0$.
\end{proof}

The trace bound gives control of positive currents in the weak topology. It does not give a bound for $\int_X(\omega+\ddc u)^n$: a bound on the integral of a trace does not bound the integral of its $n$th power.

\Needspace{19\baselineskip}
The pointwise arithmetic--geometric mean inequality relates this conserved trace to the density in a smooth Monge--Amp\`ere equation. Fix a smooth positive volume form $\mathrm{d}V$, and write $\theta_u=\theta+\ddc u$.

\begin{proposition}\label{prop:gauduchon-amgm}
For any two Hermitian forms $\eta,\gamma$,
\begin{equation}\label{eq:gauduchon-amgm}
 \eta\wedge\gamma^{n-1}
 \geq\left(\frac{\eta^n}{\gamma^n}\right)^{1/n}\gamma^n
\end{equation}
pointwise. Consequently, if a smooth potential $u$ solves
$\theta_u^n=c f\,\mathrm{d}V$ with $\theta_u>0$, $c>0$, and $f>0$, then
\begin{equation}\label{eq:gauduchon-c-bound}
 c^{1/n}
 \int_X\left(\frac{f\,\mathrm{d}V}{\gamma^n}\right)^{1/n}\gamma^n
 \leq\int_X\theta\wedge\gamma^{n-1}
\end{equation}
for every Gauduchon metric $\gamma$.
\end{proposition}
\begin{proof}
In a frame where $\gamma$ is the identity, let $\lambda_1,\ldots,\lambda_n>0$ be the eigenvalues of $\eta$. The two ratios in \eqref{eq:gauduchon-amgm} are respectively
$n^{-1}\sum_j\lambda_j$ and $(\prod_j\lambda_j)^{1/n}$. Their comparison is the arithmetic--geometric mean inequality.

Apply it with $\eta=\theta_u$, integrate, and use \eqref{eq:gauduchon-trace}.
\end{proof}

The last inequality bounds the normalization constant; it is not a top-degree mass identity. For a Hermitian background, \cref{lem:unbounded-gauduchon-trace-density} proves the corresponding measure inequality for every genuine potential by bounded mixed products and convergence of the linear currents. At a smooth positive maximizing test, its equality case gives the rigidity statement in \cref{cor:smooth-dual-threshold-rigidity}.

\section{The Gauduchon criterion for positivity}

The trace identity gives a necessary condition for $\PSH(X,\theta)$ to be nonempty: every Gauduchon degree of $\theta$ must be nonnegative. It is also sufficient. Separation produces a positive current with these prescribed pairings, and a scalar elliptic equation shows that its difference from $\theta$ is $\ddc$-exact.

\begin{theorem}[The generalized Lamari criterion]
\label{thm:lamari}
For a smooth real $(1,1)$-form $\theta$, the following conditions are equivalent:
\begin{enumerate}
\item There exists $u\in\PSH(X,\theta)$.
\item $\int_X\theta\wedge\gamma^{n-1}\geq0$ for every Gauduchon metric $\gamma$.
\end{enumerate}
\end{theorem}

\Needspace{11\baselineskip}
To apply separation, we express the test in terms of smooth strictly positive $\ddc$-closed $(n-1,n-1)$-forms $\Gamma$. Every strictly positive $(n-1,n-1)$-form has a unique strictly positive $(1,1)$-root $\Gamma=\gamma^{n-1}$. In a diagonal frame, after using the same wedge normalization for both sides, the coefficients $g_i$ of $\Gamma$ are the products $\prod_{j\neq i}\lambda_j$ of the eigenvalues of $\gamma$. Hence
\[
 \lambda_i=\frac{(\prod_jg_j)^{1/(n-1)}}{g_i}.
\]
The spectral formula uses only a positive determinant and a matrix inverse, so the root depends smoothly on the coefficients, including at repeated eigenvalues. Pointwise uniqueness makes the local roots agree on overlaps. Thus $\ddc\Gamma=0$ says exactly that this root is Gauduchon.

\begin{proof}[of \cref{thm:lamari}]
Necessity is \cref{prop:gauduchon-trace}. For sufficiency, let $F$ be the real vector space of smooth real $(n-1,n-1)$-forms and let
\[
 E=\{\Gamma\in F:\ddc\Gamma=0\},\qquad
 \ell(\Gamma)=\int_X\theta\wedge\Gamma\quad(\Gamma\in E).
\]
Fix a strictly positive $\Gamma_0\in E$, supplied by \cref{thm:gauduchon}. The hypothesis says that $\ell$ is nonnegative on strictly positive elements of $E$. By adding $\epsilon\Gamma_0$ and letting $\epsilon\downarrow0$, the same holds on semipositive elements of $E$.

For $\Phi\in F$ put
\[
 p(\Phi)=\inf\{\ell(\Gamma):\Gamma\in E,\ \Gamma\geq\Phi\}.
\]
This is a finite sublinear function. To see finiteness, choose $C$ with $-C\Gamma_0\leq\Phi\leq C\Gamma_0$. The admissible element $C\Gamma_0$ gives an upper bound, while any admissible $\Gamma$ satisfies $\ell(\Gamma)\geq-C\ell(\Gamma_0)$. Addition and positive scaling of competitors prove sublinearity. For $\Phi\in E$, positivity of $\ell(\Gamma-\Phi)$ also shows that $p(\Phi)=\ell(\Phi)$.

Hahn--Banach gives a real linear extension $T$ of $\ell$ to $F$ with $T(\Phi)\leq p(\Phi)$. If $\Phi\geq0$, the zero form is admissible for $p(-\Phi)$, so $T(\Phi)\geq0$. Moreover, the bounds by $C\Gamma_0$ show that $T$ is continuous for the uniform norm. It is therefore a positive $(1,1)$-current. The current $S=T-\theta$ annihilates $E$.

We explain why this last assertion gives $S=\ddc u$ for a distribution $u$, without using a global $\partial\bar\partial$-lemma. Fix a smooth volume form $\mathrm{d}V$ and consider
\[
 Dv=\frac{\ddc(v\Gamma_0)}{\mathrm{d}V}.
\]
This is an elliptic scalar operator with no zeroth-order term, since $\ddc\Gamma_0=0$. Its kernel consists of the constants and its index is zero. Stokes' theorem gives $\int_XDv\,\mathrm{d}V=0$. Its range is consequently exactly the smooth functions of integral zero. The inverse on functions normalized by $\int_Xv\,\mathrm{d}V=0$ is continuous in the smooth topology, by elliptic estimates.

\Needspace{9\baselineskip}
For a smooth top-degree form $\zeta$ with $\int_X\zeta=0$, solve $Dv=\zeta/\mathrm{d}V$ with this normalization and set $R\zeta=v\Gamma_0$. Then $\ddc R\zeta=\zeta$ and $R$ is continuous. Define a distribution on such top-degree forms by
\[
 \langle u,\zeta\rangle=S(R\zeta),
\]
and extend it to all top-degree forms by assigning any value on one fixed volume form of integral one. For every $\Phi\in F$, the form $\Phi-R(\ddc\Phi)$ belongs to $E$. It follows that
\[
 \langle\ddc u,\Phi\rangle
 =S(R(\ddc\Phi))=S(\Phi).
\]
Thus $\theta+\ddc u=T\geq0$. The standard local characterization of plurisubharmonic distributions, after adding a smooth strictly plurisubharmonic correction, gives the canonical quasi-plurisubharmonic representative of $u$. This representative is in $\PSH(X,\theta)$.
\end{proof}

Applying the criterion to $\theta-a\omega$ determines the largest lower bound attainable by a positive current, where $\omega$ is a fixed Hermitian form and $a\in\mathbb R$.

\Needspace{14\baselineskip}
\begin{corollary}\label{cor:gauduchon-threshold}
For a smooth real $(1,1)$-form $\theta$ and a Hermitian form $\omega$, set
\[
 \kappa_\omega(\theta)=
 \inf_{\gamma\ \mathrm{Gauduchon}}
 \frac{\int_X\theta\wedge\gamma^{n-1}}
      {\int_X\omega\wedge\gamma^{n-1}}.
\]
This is a finite real number and equals the largest real $a$ for which there is a quasi-plurisubharmonic $u$ satisfying
\[
 \theta+\ddc u\geq a\omega.
\]
The largest value is attained. In particular, $\theta$ is big precisely when $\kappa_\omega(\theta)>0$.
\end{corollary}
\begin{proof}
Smoothness gives $-C\omega\leq\theta\leq C\omega$, so the infimum lies in $[-C,C]$. The trace identity shows that any admissible $a$ is at most every displayed ratio. Conversely, with $a=\kappa_\omega(\theta)$, all Gauduchon integrals of $\theta-a\omega$ are nonnegative. Apply \cref{thm:lamari} to this smooth background. The last assertion is the definition of bigness as existence of a strictly positive representative.
\end{proof}

Vanishing of one Gauduchon degree forces a positive current to vanish.

\begin{lemma}\label{lem:zero-gauduchon-degree}
\index{Gauduchon degree!vanishing}
Let $\theta$ be a smooth real $(1,1)$-form with $\PSH(X,\theta)\ne\varnothing$.
If
\[
 \int_X\theta\wedge\gamma^{n-1}=0
\]
for one Gauduchon metric $\gamma$, then $\theta=\ddc f$ for a smooth real
function $f$. Conversely, a smoothly $\ddc$-exact form has zero degree
against every Gauduchon metric. The conclusion also holds for $n=1$, with $\gamma^{n-1}=1$.
\end{lemma}
\Needspace{8\baselineskip}
\begin{proof}
Choose $u\in\PSH(X,\theta)$ and put $S=\theta+\ddc u\geq0$. Stokes' theorem gives $\int_XS\wedge\gamma^{n-1}=0$. Since the trace measure controls all coefficient measures of $S$, one has $S=0$. Taking the trace of $\ddc u=-\theta$ with respect to $\omega$ gives an elliptic equation with smooth right-hand side. Elliptic regularity makes $u$ smooth, so $f=-u$ works. The converse is Stokes' theorem. For $n=1$, the trace measure is $S$ itself and the same argument applies.
\end{proof}

For a smooth real $(1,1)$-form $\theta$, put $a=\kappa_\omega(\theta)$. By \cref{cor:gauduchon-threshold},
$\theta-a\omega$ admits a positive current. Hence
\cref{lem:zero-gauduchon-degree} shows that the infimum defining
$\kappa_\omega(\theta)$ is attained by a smooth Gauduchon metric if and
only if
\[
 \theta-a\omega=\ddc f\quad\text{for some }f\in C^\infty(X,\mathbb R).
\]
In that case every Gauduchon metric attains it. The positive current at
the threshold always exists; a smooth minimizing test is much more restrictive.

\section{Compactness and duality of Gauduchon currents}
\label{sec:gauduchon-linear-duality}

The weak closure of normalized smooth Gauduchon forms supplies a minimizing test even when no smooth form attains the threshold. The normalization gives compactness by bounding the coefficient measures. Let $\mathscr G_\omega$ be the set of positive $(n-1,n-1)$-currents $T$ on $X$ satisfying
\index{Gauduchon current!normalized tests}
\[
 \ddc T=0,\qquad\int_X\omega\wedge T=1.
\]
These are the normalized Gauduchon currents. Let $\mathscr C_\omega$
be the weak closure of the smooth strictly positive members of
\index{Gauduchon current!closure of smooth tests}
$\mathscr G_\omega$. The root calculation for \cref{thm:lamari} identifies those members with normalized powers of
Gauduchon metrics.

\begin{proposition}\label{prop:gauduchon-linear-compact-dual}
The sets $\mathscr G_\omega$ and $\mathscr C_\omega$ are nonempty,
compact and convex. For every smooth real $(1,1)$-form $\theta$,
\begin{equation}\label{eq:gauduchon-linear-minimum}
 \kappa_\omega(\theta)
   =\min_{T\in\mathscr C_\omega}\int_X\theta\wedge T.
\end{equation}
Moreover, for $T\in\mathscr G_\omega$, membership in
$\mathscr C_\omega$ is equivalent to
\begin{equation}\label{eq:gauduchon-linear-polar}
 \int_X\theta\wedge T\geq0
 \quad\text{for every smooth real }\theta
       \text{ with }\PSH(X,\theta)\ne\varnothing.
\end{equation}
\end{proposition}
\begin{proof}
Gauduchon's theorem supplies a smooth positive member after normalization.
The fixed trace mass bounds the total variation of all coefficient
measures in finitely many local frames. Weak compactness of finite
measures therefore gives compactness of the family of positive currents
with this trace. Positivity, the trace normalization and $\ddc T=0$
pass to limits. On this bounded set, testing against smooth forms or
against continuous coefficient forms gives the same topology, by
uniform approximation. It is metrizable since the continuous test
space is separable. Thus the compactness can also be used sequentially.
The defining conditions of $\mathscr G_\omega$ are convex. Its smooth
strictly positive members form a convex set, so their closure
$\mathscr C_\omega$ is compact and convex as well.

Pairing with a fixed smooth $\theta$ is continuous. Its minimum on the
closure of the smooth tests is therefore their infimum, proving
\eqref{eq:gauduchon-linear-minimum}. The Lamari criterion shows that
\eqref{eq:gauduchon-linear-polar} holds for every smooth positive test,
and hence for every member of its weak closure.

Conversely, if $T\notin\mathscr C_\omega$, separate this point from the
compact convex set $\mathscr C_\omega$ in the weak space of currents.
A continuous linear functional for that weak topology is evaluation
against a smooth form. After changing its sign if necessary, there is
a smooth real $(1,1)$-form $\beta$ such that
\[
 \int_X\beta\wedge T
 <\min_{R\in\mathscr C_\omega}\int_X\beta\wedge R
 =\kappa_\omega(\beta).
\]
By \cref{cor:gauduchon-threshold}, the smooth form
$\beta-\kappa_\omega(\beta)\omega$ admits a positive current.
Its pairing with $T$ is negative because $\int_X\omega\wedge T=1$.
This contradicts \eqref{eq:gauduchon-linear-polar} and proves the converse.
\end{proof}

\begin{corollary}\label{cor:gauduchon-threshold-continuity}
The function $\kappa_\omega$ is concave, positively homogeneous and
invariant under addition of $\ddc f$ for $f\in C^\infty(X,\mathbb R)$. Let $\theta$ and $\eta$ be smooth real $(1,1)$-forms and let $\varepsilon\geq0$. If
$-\varepsilon\omega\leq\theta-\eta\leq\varepsilon\omega$, then
\[
 |\kappa_\omega(\theta)-\kappa_\omega(\eta)|\leq\varepsilon.
\]
If smooth real $(1,1)$-forms $\theta_j$ converge uniformly to a smooth form $\theta$
and $T_j\in\mathscr C_\omega$ minimize their pairings, every weak
cluster point of $(T_j)$ minimizes the pairing with $\theta$.
\end{corollary}
\begin{proof}
Formula \eqref{eq:gauduchon-linear-minimum} is an infimum of linear
functionals, which proves concavity and positive homogeneity.
Gauge invariance follows from $\ddc T=0$. The normalized trace gives
\[
 \left|\int_X(\theta-\eta)\wedge T\right|\leq\varepsilon
 \quad(T\in\mathscr C_\omega),
\]
and taking infima proves the estimate. For the last assertion, pass to
a subsequence $T_j\rightharpoonup T\in\mathscr C_\omega$.
The estimate and uniform convergence of the backgrounds give
\[
 \int_X\theta\wedge T
 =\lim_j\int_X\theta_j\wedge T_j
 =\lim_j\kappa_\omega(\theta_j)=\kappa_\omega(\theta).
\]
Thus $T$ is a minimizer.
\end{proof}

The closure $\mathscr C_\omega$ cannot in general be replaced by all of
$\mathscr G_\omega$. The exceptional-curve example in
\cref{ex:gauduchon-current-smoothing-obstruction} separates the two
sets by a smooth closed form. Testing against all of $\mathscr G_\omega$
will instead detect smooth positivity and nefness in
\cref{thm:smooth-cone-current-duality}.


\section{The adjoint measure of the linearization}

The Gauduchon multiplier identifies which functions lie in the range of the linearized Monge--Amp\`ere operator. Fix $u\in\HH_\omega$ and write $\eta=\omega_u$. For a smooth real direction $v$, logarithmic differentiation gives
\begin{equation}\label{eq:linearized-ma}
 L_u v=n\,\frac{\ddc v\wedge\eta^{n-1}}{\eta^n}.
\end{equation}
This is a real elliptic operator whose kernel consists of the constants. Choose the Gauduchon multiplier $\rho_u$ from \cref{thm:gauduchon}, normalized by
\[
 \int_X\rho_u\eta^n=1,
 \qquad \mathrm{d}\mu_u:=\rho_u\eta^n.
\]
The measure $\mu_u$ is a smooth positive probability measure.

\begin{proposition}\label{prop:adjoint-gauduchon}
For every smooth real $v$,
\begin{equation}\label{eq:adjoint-identity}
 \int_X L_uv\,\mathrm{d}\mu_u=0.
\end{equation}
Moreover, the equation $L_uv=h$, for smooth real $h$, has a smooth solution if and only if $\int_Xh\,\mathrm{d}\mu_u=0$. When it exists, the solution is unique modulo constants.
\end{proposition}
\Needspace{8\baselineskip}
\begin{proof}
The integral in \eqref{eq:adjoint-identity} is
\[
 n\int_X\ddc v\wedge\rho_u\eta^{n-1}
 =n\int_Xv\,\ddc(\rho_u\eta^{n-1})=0.
\]
This proves the necessary compatibility condition.

Apply the Fredholm argument from \cref{thm:gauduchon} to $\eta$. For $k\geq0$ and $0<\alpha<1$, on $C^{k+2,\alpha}(X)\to C^{k,\alpha}(X)$, the operator $L_u$ has index zero and kernel $\R$, so its range has codimension one. The nonzero functional $h\mapsto\int_Xh\,\mathrm{d}\mu_u$ annihilates that range; hence its kernel is precisely the range. This proves sufficiency. Elliptic regularity gives a smooth solution, unique modulo constants.
\end{proof}

\Needspace{10\baselineskip}
The measure $\mu_u$ depends on the solution metric $\eta$. Its multiplier is constant precisely when $\eta$ is Gauduchon: a constant multiplier satisfies the defining equation exactly in that case, and uniqueness then applies. Thus
\[
 \mu_u=\frac{\eta^n}{\int_X\eta^n}
 \quad\text{if }\ddc(\eta^{n-1})=0.
\]
On a surface, $\ddc\omega_u=\ddc\omega$, so every positive potential perturbation of a Gauduchon form is still Gauduchon. In higher dimension this persistence requires more than the Gauduchon condition at $u=0$.

\begin{proposition}\label{prop:gauduchon-potential-persistence}
Every form $\omega_u$, for $u\in\HH_\omega$, is Gauduchon if and only if
\[
 \int_X\omega_u^n=\int_X\omega^n
 \quad\text{for every }u\in\HH_\omega.
\]
In that case, with $M=\int_X\omega^n$, the normalized multiplier is $\rho_u=M^{-1}$ and $\mu_u=M^{-1}\omega_u^n$ for every $u\in\HH_\omega$.
\end{proposition}
\begin{proof}
The first variation in \cref{prop:mass-variation} is
\[
 \left.\frac{\mathrm d}{\mathrm dt}\right|_{t=0}
       \int_X\omega_{u+tv}^n
 =n\int_Xv\ddc(\omega_u^{n-1}).
\]
If all $\omega_u$ are Gauduchon, this derivative vanishes in every direction; integrating along the segment from $0$ to $u$ gives the asserted mass identity. Conversely, the identity makes the displayed integral zero for every smooth real test function $v$, so $\ddc(\omega_u^{n-1})=0$. The normalized multiplier is then the constant $M^{-1}$ by \cref{thm:gauduchon}.
\end{proof}

The differential characterization is given by \cref{thm:mass-invariance-characterization}: persistence is equivalent to $\ddc(\omega^j)=0$ for all integers $1\leq j\leq n-1$. The following example distinguishes it from the single Gauduchon identity in dimension three.

\begin{example}\label{ex:gauduchon-persistence-torus}
Let $E=\mathbb C/(\mathbb Z+\mathrm i\mathbb Z)$ and $X=E^3$. Write $z_j=x_j+\mathrm i y_j$ for the coordinates on the covering space, for $1\leq j\leq3$, and let $\eta_j=(\mathrm i/2)\mathrm dz_j\wedge\mathrm d\bar z_j$ be the descended forms. The periodic functions
\[
 b=\exp(\cos(2\pi x_1)),\qquad c=b^{-1}
\]
define a Hermitian form
\[
 \omega=\eta_1+b\eta_2+c\eta_3.
\]
Since $bc=1$,
\[
 \omega^2=2b\eta_1\wedge\eta_2
          +2c\eta_1\wedge\eta_3+2\eta_2\wedge\eta_3.
\]
The derivatives of $b,c$ involve only $x_1$, so $\ddc(\omega^2)=0$.

Put $v=\cos(2\pi x_2)$ and $h=-\pi\cos(2\pi x_2)$, so that $\ddc v=h\eta_2$. For a real $s$ with $0<|s|<1/(\pi\mathrm e)$, the form
\[
 \omega_{sv}=\eta_1+(b+sh)\eta_2+c\eta_3
\]
is positive, since $b\geq\mathrm e^{-1}$ and $|h|\leq\pi$. However,
\[
 \ddc(\omega_{sv}^2)
 =2s\ddc\omega\wedge\ddc v
 =2sh\ddc c\wedge\eta_2\wedge\eta_3\ne0.
\]
Indeed, $\ddc c=(c''/(4\pi))\eta_1$, where the prime differentiates in $x_1$, and $c''(0)=4\pi^2\mathrm e^{-1}$ while $h=-\pi$ at $x_2=0$. Thus a Gauduchon form can have arbitrarily small positive potential perturbations which are not Gauduchon.
\end{example}

Even when persistence holds, the one-form $v\mapsto\int_Xv\,\mathrm d\mu_u$ need not have a primitive on $\HH_\omega$. On a compact non-K\"ahler surface equipped with a Gauduchon form, the proposition identifies it with the raw Monge--Amp\`ere one-form divided by the constant $M$. By \cref{thm:raw-primitive}, it has no primitive. The adjoint measure provides the compatibility condition for a linear equation; this is distinct from a variational primitive.

\section{Limits of adjoint measures}
\label{sec:adjoint-measure-limits}

When all positive potential perturbations of $\omega$ are Gauduchon, their adjoint measures are normalized Monge--Amp\`ere measures. Bounded product theory then extends this formula to bounded potentials and gives continuity under convergence in capacity. The products and capacity used in this passage are constructed in \cref{prop:bounded-products,sec:hermitian-capacity}.

\begin{proposition}\label{prop:bounded-adjoint-persistence}
Assume that $\omega_v$ is Gauduchon for every $v\in\HH_\omega$, and set $M=\int_X\omega^n$. For every $u\in\PSH(X,\omega)\cap L^\infty(X)$, the measure
\[
 \mu_u:=M^{-1}(\omega+\ddc u)^n
\]
is a probability measure and agrees with the preceding definition when $u\in\HH_\omega$. If $u_j,u\in\PSH(X,\omega)\cap L^\infty(X)$, for integers $j\geq1$, are uniformly bounded and $u_j\to u$ in capacity, then $\mu_{u_j}\rightharpoonup\mu_u$.
\end{proposition}
\begin{proof}
For smooth strictly positive potentials, both assertions about the definition follow from \cref{prop:gauduchon-potential-persistence}. For bounded $u$, choose uniformly bounded smooth strictly positive approximants decreasing to $u$ by \cref{thm:bounded-exact-regularization}. Continuity of bounded products passes their common mass $M$ to the limit. Positivity is \cref{prop:bounded-products}.

To prove convergence, work on an inner coordinate ball and add one fixed smooth quadratic $q$ so that all $u_j+q$ and $u+q$ are plurisubharmonic. They are uniformly bounded and converge in local Bedford--Taylor capacity, by the capacity comparison recalled in \cref{sec:hermitian-capacity}. The classical bounded capacity-continuity theorem therefore gives convergence of their mixed products. Multiply by the fixed smooth coefficients in \eqref{eq:bounded-expansion}, then use a finite cover and a partition of unity. This gives $(\omega+\ddc u_j)^n\rightharpoonup(\omega+\ddc u)^n$ on $X$, and division by $M$ proves the assertion.
\end{proof}

The proposition applies to every Gauduchon surface. It gives an approximation-independent extension in the bounded range under its explicit geometric hypothesis. It does not construct a weak linearized operator.

For a general Hermitian form, the conformal theorem does not define $\rho_u$ when $\omega+\ddc u$ is only a current. Suppose smooth potentials $u_j\in\HH_\omega$ converge to a genuine $u\in\PSH(X,\omega)$ in capacity. Which energy or nondegeneracy hypotheses ensure that $\mu_{u_j}$ has a unique weak limit, independent of the approximation? Does this limit satisfy an adjoint identity for a weak linearization? In \cref{chap:normalization}, supporting probabilities for the normalization functional recover $\mu_u$ in the smooth case and yield a conditional approximation theorem. Outside the preceding bounded regime and those stated hypotheses, uniqueness and a distributional adjoint identity remain open here.
